% ============================================================
% CHAPITRE 6 — SPRINT 3 : FONCTIONNALITÉS ENTREPRISE ET DÉPLOIEMENT
% ============================================================

\chapter{Sprint 3 --- Fonctionnalités Entreprise et Déploiement}

\section*{Plan}

\begin{tabular}{r l r}
\textbf{1} & \textbf{Sprint Planning} & \textbf{\pageref{sec:s3_planning}} \\
\textbf{2} & \textbf{Analyse des User Stories} & \textbf{\pageref{sec:s3_stories}} \\
\textbf{3} & \textbf{Conception} & \textbf{\pageref{sec:s3_conception}} \\
\textbf{4} & \textbf{Fonctionnalités du système} & \textbf{\pageref{sec:s3_features}} \\
\textbf{5} & \textbf{Implémentation et tests} & \textbf{\pageref{sec:s3_realisation}} \\
\textbf{6} & \textbf{Sprint Rétrospective} & \textbf{\pageref{sec:s3_retro}} \\
\end{tabular}

\section*{Introduction}

Le Sprint~3 parachève le développement de la plateforme LLM Dashboard en ajoutant les fonctionnalités entreprise indispensables à son exploitation commerciale~: la facturation contractuelle avec génération automatique de factures, le traçage des sessions conversationnelles, la gestion des templates de prompts, l'administration des tenants et le déploiement conteneurisé avec intégration continue. Ce sprint transforme le prototype technique en un \textbf{produit SaaS opérationnel} prêt pour la mise en production.

% ────────────────────────────────────────────────────────────
\section{Sprint Planning}
\label{sec:s3_planning}
% ────────────────────────────────────────────────────────────

\begin{table}[H]
\centering
\caption{Sprint Backlog --- Sprint 3}
\label{tab:sprint3_backlog}
\begin{tabularx}{\textwidth}{|c|l|X|c|}
\hline
\textbf{ID} & \textbf{Module} & \textbf{User Story} & \textbf{Priorité} \\
\hline
US-12 & Billing & Gérer les contrats de facturation des tenants & Must \\
\hline
US-13 & Billing & Générer automatiquement les factures mensuelles & Must \\
\hline
US-14 & Tracing & Tracer les sessions de conversation LLM & Should \\
\hline
US-15 & Prompts & Gérer un catalogue de templates de prompts & Should \\
\hline
US-16 & Tenants & Gérer les tenants de la plateforme & Must \\
\hline
US-17 & Deploy & Déployer la plateforme via Docker et CI/CD & Must \\
\hline
\end{tabularx}
\end{table}


% ────────────────────────────────────────────────────────────
\section{Analyse des User Stories}
\label{sec:s3_stories}
% ────────────────────────────────────────────────────────────

\subsection{User Story~: Facturation contractuelle (US-12)}

\begin{table}[H]
\centering
\caption{User Story US-12~: Gestion des contrats}
\label{tab:us12}
\begin{tabularx}{\textwidth}{|l|X|}
\hline
\textbf{Titre} & Gérer les contrats de facturation des tenants \\
\hline
\textbf{En tant que} & Super Administrateur \\
\hline
\textbf{Je veux} & Créer et gérer les contrats de facturation pour chaque tenant \\
\hline
\textbf{Afin de} & Définir les conditions tarifaires adaptées à chaque client \\
\hline
\textbf{Critères d'acceptation} & \begin{minipage}[t]{\linewidth}\vspace{2pt}
\begin{itemize}[noitemsep,topsep=0pt]
\item Trois types de contrats~: \textbf{pay-as-you-go} (facturation à l'usage), \textbf{forfait} (enveloppe mensuelle fixe), \textbf{hybride} (forfait + dépassement)
\item Cycle de vie des contrats~: brouillon $\rightarrow$ proposé $\rightarrow$ accepté $\rightarrow$ actif
\item Paramètres configurables~: frais de base, seuil de tokens inclus, taux de dépassement, remise
\item Interface CRUD complète avec recherche et filtrage
\end{itemize}\vspace{2pt}
\end{minipage} \\
\hline
\end{tabularx}
\end{table}

\subsection{User Story~: Génération de factures (US-13)}

\begin{table}[H]
\centering
\caption{User Story US-13~: Génération de factures}
\label{tab:us13}
\begin{tabularx}{\textwidth}{|l|X|}
\hline
\textbf{Titre} & Générer automatiquement les factures mensuelles \\
\hline
\textbf{En tant que} & Super Administrateur \\
\hline
\textbf{Je veux} & Que le système génère automatiquement les factures en fin de mois \\
\hline
\textbf{Afin de} & Automatiser le processus de facturation et réduire les erreurs \\
\hline
\textbf{Critères d'acceptation} & \begin{minipage}[t]{\linewidth}\vspace{2pt}
\begin{itemize}[noitemsep,topsep=0pt]
\item Calcul automatique basé sur les agrégations mensuelles (table \texttt{llm\_cost\_monthly})
\item Lignes de facture détaillées par modèle LLM avec tokens et coûts unitaires
\item Support du cycle de vie~: brouillon $\rightarrow$ finalisée $\rightarrow$ envoyée $\rightarrow$ payée
\item Application automatique des conditions contractuelles (forfait, dépassement, remise)
\end{itemize}\vspace{2pt}
\end{minipage} \\
\hline
\end{tabularx}
\end{table}

\subsection{User Story~: Traçage des sessions (US-14)}

\begin{table}[H]
\centering
\caption{User Story US-14~: Traçage des sessions}
\label{tab:us14}
\begin{tabularx}{\textwidth}{|l|X|}
\hline
\textbf{Titre} & Tracer les sessions de conversation LLM \\
\hline
\textbf{En tant que} & Administrateur \\
\hline
\textbf{Je veux} & Suivre l'historique des sessions de conversation avec les LLM \\
\hline
\textbf{Afin de} & Analyser les patterns d'usage et déboguer les interactions \\
\hline
\textbf{Critères d'acceptation} & \begin{minipage}[t]{\linewidth}\vspace{2pt}
\begin{itemize}[noitemsep,topsep=0pt]
\item Corrélation parent-enfant entre les appels d'une même conversation
\item Visualisation en cascade (\textit{waterfall}) des latences
\item Suivi du nombre de tours (\textit{turns}), tokens totaux et coût par session
\item Support de l'A/B testing de prompts avec allocation de trafic
\end{itemize}\vspace{2pt}
\end{minipage} \\
\hline
\end{tabularx}
\end{table}

\subsection{User Story~: Gestion des prompts (US-15)}

\begin{table}[H]
\centering
\caption{User Story US-15~: Gestion des prompts}
\label{tab:us15}
\begin{tabularx}{\textwidth}{|l|X|}
\hline
\textbf{Titre} & Gérer un catalogue de templates de prompts \\
\hline
\textbf{En tant que} & Administrateur \\
\hline
\textbf{Je veux} & Créer, versionner et publier des templates de prompts réutilisables \\
\hline
\textbf{Afin de} & Standardiser et optimiser les prompts utilisés par l'organisation \\
\hline
\textbf{Critères d'acceptation} & \begin{minipage}[t]{\linewidth}\vspace{2pt}
\begin{itemize}[noitemsep,topsep=0pt]
\item CMS de templates avec versionnement (v1, v2, v3\ldots)
\item Support des variables Jinja2 avec rendu dynamique
\item Cycle de publication~: brouillon $\rightarrow$ publié $\rightarrow$ archivé
\item Catégorisation par tags et recherche textuelle
\end{itemize}\vspace{2pt}
\end{minipage} \\
\hline
\end{tabularx}
\end{table}


% ────────────────────────────────────────────────────────────
\section{Conception}
\label{sec:s3_conception}
% ────────────────────────────────────────────────────────────

\subsection{Diagramme de classes --- Module Billing}

Le diagramme de classes du module Billing modélise le système de facturation contractuelle avec ses trois types de contrats et le cycle de vie des factures~:

\begin{figure}[H]
    \centering
    \fbox{\parbox{14cm}{\centering\vspace{4cm}\textit{Diagramme de classes --- Module Billing}\\[0.5cm]\textit{(Classes~: BillingContract, Invoice, InvoiceLineItem, BillingService)}\\[0.5cm]\textit{(Enums~: ContractType, ContractStatus, InvoiceStatus)}\\[0.3cm]\textit{(Voir fichier PlantUML~: ch6\_class\_billing.puml)}\vspace{3cm}}}
    \caption{Diagramme de classes du module Billing}
    \label{fig:class_billing}
\end{figure}

\subsection{Diagramme de classes --- Modules Tracing et Prompts}

Le diagramme suivant modélise les entités nécessaires au traçage des sessions conversationnelles, à la gestion des templates de prompts et aux alertes multi-canal~:

\begin{figure}[H]
    \centering
    \fbox{\parbox{14cm}{\centering\vspace{4cm}\textit{Diagramme de classes --- Modules Tracing et Prompts}\\[0.5cm]\textit{(Classes~: LLMSession, LLMTokenLog, TracingService, ABTestConfig,}\\
    \textit{PromptTemplate, PromptService, AlertRule)}\\[0.5cm]\textit{(Enums~: SessionStatus, PromptStatus, AlertChannel)}\\[0.3cm]\textit{(Voir fichier PlantUML~: ch6\_class\_tracing.puml)}\vspace{3cm}}}
    \caption{Diagramme de classes des modules Tracing et Prompts}
    \label{fig:class_tracing}
\end{figure}

\subsection{Diagramme de séquence --- Traçage d'une session LLM}

Le diagramme suivant illustre le cycle de vie complet d'une session conversationnelle~: création automatique, ajout de tours successifs, fermeture et consultation de la vue en cascade~:

\begin{figure}[H]
    \centering
    \fbox{\parbox{14cm}{\centering\vspace{4cm}\textit{Diagramme de séquence --- Traçage d'une session}\\[0.5cm]\textit{(Agent $\rightarrow$ Gateway $\rightarrow$ TracingService $\rightarrow$ DB $\rightarrow$ EventBus $\rightarrow$ Dashboard)}\\[0.5cm]\textit{(1. Création session \quad 2. Ajout tours successifs \quad 3. Fermeture)}\\
    \textit{(4. Consultation waterfall avec latences et tokens par tour)}\\[0.3cm]\textit{(Voir fichier PlantUML~: ch6\_seq\_tracing.puml)}\vspace{3cm}}}
    \caption{Diagramme de séquence du traçage d'une session LLM}
    \label{fig:seq_tracing}
\end{figure}

\subsection{Diagramme de séquence --- Gestion des prompts}

Le flux de gestion des templates de prompts comprend la création, la publication, le versionnement et le rendu dynamique avec des variables Jinja2~:

\begin{figure}[H]
    \centering
    \fbox{\parbox{14cm}{\centering\vspace{4cm}\textit{Diagramme de séquence --- Gestion des templates de prompts}\\[0.5cm]\textit{(Admin $\rightarrow$ Frontend $\rightarrow$ API $\rightarrow$ PromptService $\rightarrow$ DB $\rightarrow$ Jinja2Engine)}\\[0.5cm]\textit{(1. Création DRAFT \quad 2. Publication \quad 3. Nouvelle version)}\\
    \textit{(4. Rendu dynamique avec variables)}\\[0.3cm]\textit{(Voir fichier PlantUML~: ch6\_seq\_prompt\_mgmt.puml)}\vspace{3cm}}}
    \caption{Diagramme de séquence de la gestion des prompts}
    \label{fig:seq_prompts}
\end{figure}

\subsection{Diagramme de séquence --- Facturation}

\begin{figure}[H]
    \centering
    \fbox{\parbox{14cm}{\centering\vspace{4cm}\textit{Diagramme de séquence --- Génération automatique de factures}\\[0.5cm]\textit{(Celery Beat $\rightarrow$ BillingService $\rightarrow$ ContractRepo $\rightarrow$ CostMonthlyRepo)}\\[0.5cm]\textit{(1. Récupération contrat actif \quad 2. Calcul selon type)}\\
    \textit{(3. Application remise \quad 4. Création facture avec lignes détaillées)}\\[0.3cm]\textit{(Voir fichier PlantUML~: ch6\_seq\_billing.puml)}\vspace{3cm}}}
    \caption{Diagramme de séquence de la génération de factures}
    \label{fig:seq_billing}
\end{figure}

\subsection{Diagrammes d'états --- Cycles de vie}

Les diagrammes suivants modélisent les machines à états des trois entités dont le cycle de vie est régi par des transitions formalisées~:

\begin{figure}[H]
    \centering
    \fbox{\parbox{14cm}{\centering\vspace{3cm}\textit{Diagramme d'états --- Cycle de vie d'un contrat}\\[0.5cm]\textit{(DRAFT $\rightarrow$ PROPOSED $\rightarrow$ ACCEPTED $\rightarrow$ ACTIVE $\rightarrow$ TERMINATED)}\\[0.3cm]\textit{(Voir fichier PlantUML~: ch6\_state\_contract.puml)}\vspace{2cm}}}
    \caption{Machine à états du cycle de vie d'un contrat}
    \label{fig:state_contract}
\end{figure}

\begin{figure}[H]
    \centering
    \fbox{\parbox{14cm}{\centering\vspace{3cm}\textit{Diagramme d'états --- Cycle de vie d'une facture}\\[0.5cm]\textit{(DRAFT $\rightarrow$ FINALIZED $\rightarrow$ SENT $\rightarrow$ PAID / OVERDUE)}\\[0.3cm]\textit{(Voir fichier PlantUML~: ch6\_state\_invoice.puml)}\vspace{2cm}}}
    \caption{Machine à états du cycle de vie d'une facture}
    \label{fig:state_invoice}
\end{figure}

\begin{figure}[H]
    \centering
    \fbox{\parbox{14cm}{\centering\vspace{3cm}\textit{Diagramme d'états --- Cycle de vie d'un template de prompt}\\[0.5cm]\textit{(DRAFT $\rightarrow$ PUBLISHED $\rightarrow$ ARCHIVED / NEW\_VERSION)}\\[0.3cm]\textit{(Voir fichier PlantUML~: ch6\_state\_prompt.puml)}\vspace{2cm}}}
    \caption{Machine à états du cycle de vie d'un template de prompt}
    \label{fig:state_prompt}
\end{figure}


% ────────────────────────────────────────────────────────────
\section{Fonctionnalités du système}
\label{sec:s3_features}
% ────────────────────────────────────────────────────────────

Cette section présente les principales interfaces de la plateforme LLM Dashboard, offrant une vue d'ensemble des fonctionnalités accessibles aux différents utilisateurs.

\subsection{Page d'accueil --- Tableau de bord exécutif}

\begin{figure}[H]
    \centering
    \fbox{\parbox{14cm}{\centering\vspace{4cm}\textit{Capture d'écran --- Tableau de bord exécutif}\\[0.5cm]\textit{(KPIs : Tokens totaux, Coût total, Requêtes, Taux d'erreur)}\\[0.5cm]\textit{(Graphiques de tendance, répartition par modèle, top tenants)}\vspace{3cm}}}
    \caption{Page d'accueil --- Tableau de bord exécutif}
    \label{fig:feat_dashboard}
\end{figure}

\subsection{Page Anomalies --- Détection et explication IA}

\begin{figure}[H]
    \centering
    \fbox{\parbox{14cm}{\centering\vspace{4cm}\textit{Capture d'écran --- Page des anomalies}\\[0.5cm]\textit{(Liste filtrable par sévérité, algorithme, date)}\\[0.5cm]\textit{(Détail avec explication IA, recommandations et anomalies similaires)}\vspace{3cm}}}
    \caption{Page Anomalies --- Détection et explication IA}
    \label{fig:feat_anomalies}
\end{figure}

\subsection{Page Gouvernance --- Configuration des règles}

\begin{figure}[H]
    \centering
    \fbox{\parbox{14cm}{\centering\vspace{4cm}\textit{Capture d'écran --- Page de gouvernance}\\[0.5cm]\textit{(Tableau des règles actives avec type, seuil et scope)}\\[0.5cm]\textit{(Formulaire de création : rate\_limit, budget\_cap, model\_restriction, etc.)}\vspace{3cm}}}
    \caption{Page Gouvernance --- Configuration des règles}
    \label{fig:feat_governance}
\end{figure}

\subsection{Page Facturation --- Contrats et factures}

\begin{figure}[H]
    \centering
    \fbox{\parbox{14cm}{\centering\vspace{4cm}\textit{Capture d'écran --- Page de facturation}\\[0.5cm]\textit{(Onglets : Factures / Paramètres contrat / Tarification)}\\[0.5cm]\textit{(Détail d'une facture avec lignes par modèle et montants)}\vspace{3cm}}}
    \caption{Page Facturation --- Contrats et factures}
    \label{fig:feat_billing}
\end{figure}

\subsection{Page Playground --- Test d'appels LLM}

\begin{figure}[H]
    \centering
    \fbox{\parbox{14cm}{\centering\vspace{4cm}\textit{Capture d'écran --- Playground LLM}\\[0.5cm]\textit{(Éditeur Monaco pour le prompt, sélecteur de modèle/provider)}\\[0.5cm]\textit{(Résultat avec tokens consommés, coût et latence)}\vspace{3cm}}}
    \caption{Page Playground --- Interface de test d'appels LLM}
    \label{fig:feat_playground}
\end{figure}

\subsection{Page Prévisions --- Forecasting budgétaire}

\begin{figure}[H]
    \centering
    \fbox{\parbox{14cm}{\centering\vspace{4cm}\textit{Capture d'écran --- Page de prévisions budgétaires}\\[0.5cm]\textit{(3 courbes~: optimiste, moyen, pessimiste)}\\[0.5cm]\textit{(Alertes de risque et horizon configurable)}\vspace{3cm}}}
    \caption{Page Prévisions --- Forecasting budgétaire}
    \label{fig:feat_forecasting}
\end{figure}

\subsection{Page Sessions --- Traçage des conversations}

\begin{figure}[H]
    \centering
    \fbox{\parbox{14cm}{\centering\vspace{4cm}\textit{Capture d'écran --- Page de traçage des sessions}\\[0.5cm]\textit{(Liste des sessions avec durée, tours, tokens totaux)}\\[0.5cm]\textit{(Vue waterfall des appels individuels)}\vspace{3cm}}}
    \caption{Page Sessions --- Traçage des conversations LLM}
    \label{fig:feat_sessions}
\end{figure}

\subsection{Page Tenants --- Administration multi-tenant}

\begin{figure}[H]
    \centering
    \fbox{\parbox{14cm}{\centering\vspace{4cm}\textit{Capture d'écran --- Page de gestion des tenants}\\[0.5cm]\textit{(Tableau des tenants avec statut, nombre d'utilisateurs, consommation)}\\[0.5cm]\textit{(Actions : créer, activer, désactiver, configurer le contrat)}\vspace{3cm}}}
    \caption{Page Tenants --- Administration multi-tenant}
    \label{fig:feat_tenants}
\end{figure}


% ────────────────────────────────────────────────────────────
\section{Implémentation et tests}
\label{sec:s3_realisation}
% ────────────────────────────────────────────────────────────

\subsection{Module Billing --- Implémentation}

Le module de facturation implémente un système complet de gestion des contrats et de génération de factures~:

\begin{lstlisting}[language=Python3, caption={Service de génération de factures (modules/billing/services/billing\_service.py)}]
class BillingService:
    async def generate_monthly_invoice(self, tenant_id: str,
                                        month: int, year: int) -> Invoice:
        # 1. Recuperer le contrat actif du tenant
        contract = await self.contract_repo.get_active(tenant_id)
        
        # 2. Recuperer les couts mensuels agreges
        monthly_costs = await self.cost_repo.get_monthly(
            tenant_id, month, year
        )
        
        # 3. Calculer selon le type de contrat
        if contract.contract_type == "pay_as_you_go":
            total = sum(c.total_cost for c in monthly_costs)
        elif contract.contract_type == "flat_rate":
            total = contract.base_fee
        else:  # hybrid
            usage = sum(c.total_cost for c in monthly_costs)
            if usage > contract.included_amount:
                overage = (usage - contract.included_amount) 
                          * contract.overage_rate
                total = contract.base_fee + overage
            else:
                total = contract.base_fee
        
        # 4. Appliquer la remise
        total *= (1 - contract.discount_pct / 100)
        
        # 5. Creer la facture avec les lignes detaillees
        invoice = await self.invoice_repo.create(
            tenant_id=tenant_id, total=total,
            line_items=self._build_line_items(monthly_costs),
        )
        return invoice
\end{lstlisting}


\subsection{Déploiement Docker et CI/CD}

Le déploiement de la plateforme repose sur \textbf{Docker Compose} pour l'orchestration des six services conteneurisés et sur \textbf{GitHub Actions} pour l'intégration continue~:

\begin{lstlisting}[language=bash, caption={Déploiement Docker Compose (extrait de docker-compose.yml)}]
services:
  postgres:
    image: timescale/timescaledb:latest-pg15
    healthcheck:
      test: ["CMD-SHELL", "pg_isready -U llmuser"]
  
  redis:
    image: redis:7-alpine
    command: redis-server --maxmemory 256mb
  
  api:
    build: .
    depends_on:
      postgres: { condition: service_healthy }
      redis:    { condition: service_healthy }
  
  worker:
    build: .
    command: celery -A celery_app worker --concurrency=2
  
  scheduler:
    build: .
    command: celery -A celery_app beat
  
  nginx:
    image: nginx:alpine
    ports: ["80:80", "443:443"]
\end{lstlisting}

Le pipeline CI/CD comprend quatre étapes exécutées automatiquement à chaque push~:

\begin{enumerate}[noitemsep]
    \item \textbf{Lint}~: vérification syntaxique de tous les fichiers Python~;
    \item \textbf{Test}~: exécution de pytest avec services TimescaleDB et Redis éphémères~;
    \item \textbf{Build}~: construction de l'image Docker multi-stage et push vers GHCR~;
    \item \textbf{Deploy}~: notification et déploiement sur le serveur de production.
\end{enumerate}


\subsection{Tests du Sprint~3}

\begin{table}[H]
\centering
\caption{Résultats des tests --- Sprint 3}
\label{tab:tests_s3}
\begin{tabular}{|l|c|c|c|}
\hline
\textbf{Module} & \textbf{Tests unitaires} & \textbf{Tests d'intégration} & \textbf{Résultat} \\
\hline
Billing & 14 & 6 & \textcolor{ForestGreen}{20/20 PASS} \\
\hline
Tracing & 8 & 4 & \textcolor{ForestGreen}{12/12 PASS} \\
\hline
Prompts & 7 & 3 & \textcolor{ForestGreen}{10/10 PASS} \\
\hline
Tenants & 5 & 2 & \textcolor{ForestGreen}{7/7 PASS} \\
\hline
Déploiement & -- & 4 & \textcolor{ForestGreen}{4/4 PASS} \\
\hline
\textbf{Total Sprint 3} & \textbf{34} & \textbf{19} & \textcolor{ForestGreen}{\textbf{53/53 PASS}} \\
\hline
\end{tabular}
\end{table}

\subsection{Synthèse globale des tests}

\begin{table}[H]
\centering
\caption{Synthèse globale des tests de la plateforme}
\label{tab:tests_global}
\begin{tabular}{|l|c|c|c|}
\hline
\textbf{Sprint} & \textbf{Tests unitaires} & \textbf{Tests d'intégration} & \textbf{Total} \\
\hline
Sprint 1 (Infrastructure) & 35 & 16 & 51 \\
\hline
Sprint 2 (Intelligence Artificielle) & 55 & 21 & 76 \\
\hline
Sprint 3 (Entreprise) & 34 & 19 & 53 \\
\hline
\textbf{Total plateforme} & \textbf{124} & \textbf{56} & \textcolor{ForestGreen}{\textbf{180/180 PASS}} \\
\hline
\end{tabular}
\end{table}


% ────────────────────────────────────────────────────────────
\section{Sprint Rétrospective}
\label{sec:s3_retro}
% ────────────────────────────────────────────────────────────

\begin{table}[H]
\centering
\caption{Rétrospective du Sprint 3}
\label{tab:retro_s3}
\begin{tabularx}{\textwidth}{|l|X|}
\hline
\textbf{Ce qui a bien fonctionné} & \begin{minipage}[t]{\linewidth}\vspace{2pt}
\begin{itemize}[noitemsep,topsep=0pt]
\item Le module de facturation contractuelle couvre les trois types de contrats industriels
\item Le déploiement Docker Compose permet un démarrage complet en une seule commande
\item Le pipeline CI/CD automatise l'ensemble du processus de livraison
\item L'interface de gestion des tenants offre une vue centralisée aux super administrateurs
\end{itemize}\vspace{2pt}
\end{minipage} \\
\hline
\textbf{Ce qui peut être amélioré} & \begin{minipage}[t]{\linewidth}\vspace{2pt}
\begin{itemize}[noitemsep,topsep=0pt]
\item La génération PDF des factures peut être enrichie avec un template plus élaboré
\item Le système d'alertes multi-canal (email, Slack, webhook) nécessite une configuration plus poussée
\item Les tests de charge en conditions réelles doivent être réalisés avant le déploiement à grande échelle
\end{itemize}\vspace{2pt}
\end{minipage} \\
\hline
\end{tabularx}
\end{table}


\section*{Conclusion}

Le Sprint~3 a achevé le développement de la plateforme LLM Dashboard en ajoutant les fonctionnalités entreprise essentielles à son exploitation commerciale. Le module de facturation contractuelle, avec ses trois types de contrats et sa génération automatique de factures, répond aux besoins spécifiques de TIM Tunisie pour la facturation de ses clients QALITAS et GMAO PRO. Le traçage des sessions, la gestion des prompts et l'administration des tenants complètent l'offre fonctionnelle. Le déploiement conteneurisé via Docker Compose et le pipeline CI/CD GitHub Actions garantissent une mise en production reproductible et automatisée. Au total, la plateforme est validée par \textbf{180 tests} couvrant l'ensemble des 11 modules fonctionnels.
